\chapter{Гейт единой карты весов}

После провала простого ярусного правила полезно отделить три различных
утверждения: существование числовой таблицы весов, происхождение этой
таблицы из одного действия и логическую независимость оставшихся
динамических разрывов. Эти утверждения не эквивалентны.

\section{Фиксация опорной метрики}

Карта весов \(W\) определена только относительно заранее фиксированного
скалярного произведения. Поэтому варианты ``\(W=1\) с каноническими
нормами'' и ``\(W=1\) с сырыми нормами'' не являются двумя правилами одного
класса: второй вариант уже заменяет базовую метрику. В каноническом базисе
целевые значения равны
\begin{align}
 \rho_\tau&=\pi^2+2\pi+\frac23,\label{eq:WgateTau20260807}\\
 \rho_\nu&=23+\pi^{-1}.\label{eq:WgateNu20260807}
\end{align}

Для постоянного веса \(W=1\) получаются
\begin{equation}
 \rho_\tau^{C0}=\frac83,
 \qquad
 \rho_\nu^{C0}=24,
\end{equation}
так что класс \(C0\) проваливает оба целевых значения.

\section{Классы \(C1\) и \(C2\)}

Правило, использующее только степень формы и наличие квантованного периода,
назначает обычным нуль-формам единичный вес, а интегральной одной форме ---
вес \(\pi^{-1}\). Оно даёт
\begin{equation}
 \rho_\tau^{C1}=\frac83,
 \qquad
 \rho_\nu^{C1}=23+\pi^{-1}.
\end{equation}
Следовательно, \(C1\) проходит нейтринный сектор и проваливает лептонный.

Формулировка ``все нуль-формы имеют сырую норму'' не принадлежит \(C1\):
сырая норма равна \(\pi^2\) на \(\mathbb{RP}^3\), \(2\pi\) на \(S^1\) и
другому числу на ином факторе. Это уже зависимость от фактора, то есть
класс \(C2\).

В \(C2\) действительно существует таблица
\begin{equation}
 W_{(\mathbb{RP}^3,0)}=\pi^2,
 \quad
 W_{(S^1,0)}=2\pi,
 \quad
 W_{(\mathrm{fin},0)}=1,
 \quad
 W_{(S^1,1)_{\mathrm{int}}}=\pi^{-1},
\end{equation}
которая одновременно воспроизводит \eqref{eq:WgateTau20260807} и
\eqref{eq:WgateNu20260807}.
Она не использует названия ``лептон'' и ``нейтрино'', поэтому формально не
является запрещённым правилом \(C3\). Однако разбиение на четыре ячейки уже
содержит все требуемые относительные нормировки. Без вариационного
функционала, симметрии или граничной симплектической формы это согласованная
таблица метрики, но не предсказание.

\begin{nogocriterion}[Табличное прохождение не равно единой мере]
Прохождение двух агрегатных норм таблицей \(C2\) не считается прохождением
двухсекторного физического гейта, пока:
\begin{enumerate}
 \item веса ячеек не получены как гессиан одного заранее заданного действия;
 \item не задан закон их совместного преобразования с вершинами при смене
 координат поля;
 \item тем же правилом не вычислена независимая наблюдаемая, не
 использованная при построении таблицы.
\end{enumerate}
\end{nogocriterion}

\section{Гейт ковариантности и логический статус}

При замене координаты поля \(x'=a x\) числовой коэффициент квадратичной
формы меняется обратно квадрату масштаба, а коэффициенты вершин должны
преобразоваться согласованно. Поэтому физическим объектом является не
список норм сам по себе, а ковариантная пара ``метрика плюс карта
наблюдаемой''. Предложенная таблица \(C2\) такого закона пока не задаёт.

Провал \(C0\)--\(C1\) и недовыведенность \(C2\) показывают, что единая мера
в текущем формализме не построена. Они не доказывают, что более глубокое
действие не может вывести \(C2\), и не доказывают логическую независимость
разрывов меры, CP-источника, выбора вакуума и порогов.

\begin{proposition}[Вердикт \(W\)-гейта]
Алгебраический статус равен
\begin{equation}
 C0:\ \mathrm{fail},
 \qquad
 C1:\ \mathrm{one\ sector},
 \qquad
 C2:\ \mathrm{lookup\ pass}.
\end{equation}
Операторный и двухсекторный предсказательный гейты для \(C2\) провалены,
поскольку родительское происхождение и ковариантность весов отсутствуют.
Число новых закрытых физических предсказаний остаётся нулевым.
\end{proposition}

\begin{protocol}[Условие переоткрытия \(W\)-гейта]
Следующая модель должна вывести указанную \(C2\)-метрику как гессиан или
граничную форму одного действия до обращения к массам, задать совместное
преобразование stiffness и couplings и затем пройти независимый
электромагнитный либо роторный тест. Только после этого допустимо проверять,
являются ли остальные динамические разрывы проекциями одного принципа.
\end{protocol}

Машинный протокол сохранён в
\texttt{s2t\_weight\_map\_gate\_results.json}.